\subsection*{Justificacion Fila 7:}

La operación indica tomar la columna \texttt{R.B}, y luego restarle las celdas de la columna \texttt{S.B} tales que su respectivo registro \texttt{S.C} sea mayor o igual a 3.
\begin{center}
\begin{tabular}{c c c c c c c c}
    $\pi_B(R) =$
&
\begin{tabular}{|c|}
\hline
\textbf{B} \\ \hline
  x
  \\
  y
  \\
  z
  \\
  x
  \\
  a 
  \\ \hline
\end{tabular}
& &     $\pi_B(\sigma_{C \ge 3}(S)) =$ &
\begin{tabular}{|c|}
\hline
\textbf{B} \\ \hline
  y
  \\
  w
  \\
  y
  \\ \hline
\end{tabular}
& &     $\pi_B(R)-\pi_B(\sigma_{C \ge 3}(S)) =$ &
\begin{tabular}{|c|}
\hline
\textbf{B} \\ \hline
  x
  \\
  z
  \\
  a
  \\ \hline
\end{tabular}
\end{tabular}
\end{center}

Por lo tanto, la cardinalidad del resultado es \textbf{3}.